\section{A general differentiable forward-wall net}
\label{sec:ad}
\subsection{Differentiate the solved analysis, including its boundaries}
The requested object is a prescribed-wall flow operator with wall,
inlet/gas, symmetry and exit dependence. It must include the starting
line, every propagated parent, wall intersections/tangency, axis
reflections, post-D/BDE regions where used, and the final momentum,
pressure, mass and energy integrals.

\current The host bounded forward net and JAX throat kernel provide
useful references. The JAX \code{build_start_line} path still raises
\code{NotImplementedError}. The sensitivity API supplies selected
fixed-state partials and explicitly defers total solved-design and
bell-wall derivatives; it refuses the exact-D mode used by the recent
accepted profile. This is not an end-to-end differentiable arbitrary-wall
analysis or a total-gradient API for that profile.

\proposal Reuse the host implementation as a primal oracle within its
established domain, while separating continuous unit-process solves
from discrete topology construction. The first full differentiable
target should be attached, smooth, calorically perfect flow with
matched starting data. Add other physical regimes after this complete
chain is validated.

\subsection{Record the topology and its validity neighborhood}
A proposed topology record should include parent indices, row masks and
widths, raw/special insertions, wall and axis reflection roles, interpolation
brackets, intersection branches, shock-region labels, exit cropping and
termination reasons. Within one accepted topology, solve continuous
unit-process residuals with explicit convergence and conditioning checks.

The presence of \code{stop_gradient} or a fixed-shape mask does not prove
that a discrete decision is irrelevant to the physical derivative.
Changing a wall-intersection bracket, reflection count, active constraint,
shock attachment, separation pattern, mesh topology or property branch
can change the differentiated problem. A physically continuous output
may still have a non-smooth discrete approximation or different one-sided
derivatives.

A practical proposed policy is to diagnose an admissible local topology
neighborhood, use trust-region steps within it, and rebuild/revalidate
when events occur. At an event, mark the ordinary smooth derivative
unavailable or provide explicitly validated one-sided/branch derivatives.
Smoothing a switch is a modeling/numerical approximation whose bias must
be measured; it is not automatically the derivative of the unsmoothed
flow.

\subsection{Root, optimization and iteration derivatives are different}
For a converged square residual
\begin{equation}
F(z,x)=0,\qquad
F_z z_x=-F_x,
\label{eq:IFT}
\end{equation}
implicit sensitivities require a differentiable local residual and
nonsingular state Jacobian. A transposed linear solve gives an adjoint
form for scalar output \(J(z,x)\):
\begin{equation}
F_z^\mathsf{T}\lambda=J_z^\mathsf{T},
\qquad
\frac{\dd J}{\dd x}=J_x-\lambda^\mathsf{T}F_x.
\end{equation}
Blondel et al.\ give the root/optimality-condition framework and discuss
approximate-solve error \cite[\S2.1, PDF pp.~3--4;
\S3, PDF p.~6]{blondel}. These hypotheses are requirements to check for
the chosen MOC system, not properties conferred by using an AD library.

For an overdetermined least-squares calculation with a nonzero final
residual, the stationarity or active-set/KKT system is the relevant
implicit object. For fixed symmetric weights \(W\),
\begin{align}
\phi(z,x)&=\tfrac12r^\mathsf{T}Wr,&
g(z,x)&=J_r^\mathsf{T}Wr=0,\\
g_z&=J_r^\mathsf{T}WJ_r+
\sum_i(Wr)_i\nabla_z^2r_i.
\end{align}
This displayed Hessian follows by differentiating that objective.
Dropping its residual-curvature term is a Gauss--Newton approximation,
not automatically the derivative of the exact optimum. Reduced D/E
parameterizations, design-dependent weights/bounds and stable active
constraints require their full chain rules.

Differentiating a fixed number of nonlinear iterations gives the
derivative of that algorithmic output. It approaches a solved-state
derivative only under additional convergence/regularity conditions.
The existing library supports multiple adjoint strategies, but installed
behavior and the actual differentiated residual must be audited before
selecting one \cite{optimistix}. No present implementation defect in
the library is asserted here.

\subsection{What a meaningful derivative test looks like}
\begin{enumerate}
\item Check the full primal result against the host oracle on the same
wall, gas state, topology and numerical policy.
\item Verify each continuous gas, intersection, wall, axis and exit
operator, including implicit-root contributions.
\item Perform full-output directional Taylor tests:
\[
E(h)=\left|J(x+hv)-J(x)-h\,D J(x)v\right|.
\]
For a smooth sufficiently accurate solve, the first-order remainder
should show the expected second-order regime before roundoff/primal
error dominates. Sweep step sizes; one finite-difference agreement
does not establish a useful error regime.
\item Repeat under physical-grid, property-grid and nonlinear-tolerance
refinement. Primal error must be below the tested signal, and adjoint
residual/conditioning should be recorded.
\item At topology or physical-regime events, compare admissible one-sided
branches rather than treating an across-event centered difference as a
smooth-gradient oracle.
\end{enumerate}
Test wall controls, \(\Rt,R_u,R_d,\epsilon,L\), chamber state and gas
parameters wherever supported. For engine totals include chamber
thermochemistry, composition, starting-line, inner wall temperature,
throat/mass-flow closure and the output's direct parameter dependence.

\subsection{Cost and memory are part of the design}
An adjoint can avoid forming every design-state Jacobian column, but it
does not make the primal analysis free. Padded characteristic histories,
nonlinear correctors, equilibrium/species solves and reverse-mode storage
can all be large. The current throat representation also grows as
rows widen. Measure scaling with actual accepted problem sizes.

Sparse/local residual structure, converged unit-process implicit solves,
checkpointing, branch-consistent warm starts and compiled reuse are
candidate techniques. Cache keys must bind geometry, gas/data identity,
boundary conditions, topology and numerical policy. Reusing a previous
converged root is an initialization strategy; accepting an unconverged
history-dependent result would change the analysis.

The forward-wall operator is a useful priority because it can avoid
solving a fresh inverse contour-optimization problem for each outer wall
candidate. Any cost advantage remains to be demonstrated by profiling.
The recently measured 25-minute pair cannot be assumed suitable for
thousands of nested MDO evaluations.

\section{Effects on full-engine MDO}
\label{sec:mdo}
\subsection{The state closure must use the same aerodynamics as the objective}
\proposal A coupled formulation can collect throat/mass-flow states
\(y=(\Rt,\dot m)\), flow/BL states \(z\), wall/coolant states \(w\), and
design \(x\) into residuals of the form
\begin{align}
R_F&=F_{\mathrm{flow}}(z,w,y,x;p_a)-F_{\mathrm{req}},\\
R_m&=\dot m-A_tG_{*,\mathrm{del}}(y,x;\mathcal G),\\
R_{\mathrm{flow}}&=0,\qquad
R_{\mathrm{BL/thermal}}=0.
\end{align}
This is an architecture proposal. Here \(G_{*,\mathrm{del}}\) is the
mass flux averaged over the declared reference throat area; a
nonuniform throat requires the corresponding mass-flow integral.
The delivered sonic mass-flux model
must declare combustion/loss conventions rather than double-counting
an existing \(c^*\) efficiency. For each regime, performance should come
from a declared consistent control volume with momentum, pressure,
mass and energy accounting.

Whether implemented monolithically or by nested solves, the linearized
coupling must be present in the total derivative. For example,
\begin{equation}
\begin{bmatrix}
R_y & R_z\\ F_y & F_z
\end{bmatrix}
\begin{bmatrix}y_x\\ z_x\end{bmatrix}
=-
\begin{bmatrix}R_x\\F_x\end{bmatrix}
\end{equation}
shows why a frozen inner flow or an unchanged Q1D sizing residual loses
part of the design sensitivity. Wall/coolant states add corresponding
blocks. Block elimination, Schur complements and adjoints are numerical
choices to evaluate after the residual is physically defined.

\subsection{Propagation through the engine}
\begin{longtable}{@{}P{0.22\linewidth}P{0.34\linewidth}P{0.36\linewidth}@{}}
\toprule\textbf{Discipline} & \textbf{New nozzle information} & \textbf{MDO consequence}\\\midrule
\endhead
Thrust, throat and propellant flow &
Accepted \(F,C_F,G_*,c^*\), gas mode and pressure boundary. &
Changes \(\Rt,\dot m\), O/F split, \(I_{sp}\), and the derivatives of
the thrust/mass-flow residuals.\\
Chamber and nozzle geometry &
Physical wall, length/curvature, throat area and effective core boundary. &
Changes chamber \(L^*A_t\) volume closure, envelope, wall area, mass
and the common station grid.\\
Heat transfer and regeneration &
Edge \(p,T,\rho,V,Y\), transport, wall heat flux and cooling condition. &
Changes wall/coolant temperatures, coolant enthalpy rise, channel
sizing and thermal/strength constraints.\\
Hydraulics and pumps &
Propellant split, coolant density/viscosity, jacket path and pressure loss. &
Changes required pump rise, volume flow, efficiency, shaft power,
NPSH and tip-speed screens.\\
Electric package and battery &
Pump duty and operating duration/efficiencies. &
Changes motor/inverter/pump screening masses and energy/power
battery branches.\\
Structure and loads &
Wall and coolant pressure, temperature, effective material geometry,
separation/load evidence. &
Changes wall/liner/closeout stresses and margins; side-load
constraints need a separate supported model.\\
Objectives and admission &
Common model identity, validity, uncertainty, total derivatives and cost. &
Changes candidate feasibility and optimum; unsupported physics
must remain unknown/infeasible rather than silently accepted.\\
\bottomrule
\end{longtable}

These are coupling pathways, not predictions that every higher-fidelity
model improves performance or reduces mass. At fixed target thrust,
\[
I_{sp}=\frac{F}{\dot m g_0},\qquad
\frac{\dd I_{sp}}{I_{sp}}=
\frac{\dd F}{F}-\frac{\dd\dot m}{\dot m}.
\]
For a simple incompressible duty model,
\[
P_{\mathrm{elec}}\simeq
\frac{\dot m\,\Delta p}{\rho\,\eta_p\eta_{\mathrm{drive}}}.
\]
A nozzle change can reduce required mass flow while increasing wall
length, heat pickup or jacket pressure drop. The net mass/efficiency
tradeoff must be evaluated through the coupled model.

\subsection{Existing thermal and property seams need reconciliation}
Current MDO cooling uses mission-level gas \(c_p,\mu,\Pr\) and coolant
\(\rho,c_p,\mu,k\). A film core-velocity helper also reads the mission gas
constant while chamber surfaces may supply another active value.
These are explicit integration boundaries to audit before claiming
spatially consistent variable-gas MDO. This inspection does not quantify
their error at a particular engine.

The host real-fluid, radiation and manifold tools are useful
infrastructure, but their presence does not establish that the MDO
cooling march calls them. Hot-gas thermodynamic consistency and
real-fluid coolant modeling are distinct workstreams. Film composition,
wall chemistry and heat-loss coupling also require explicit assumptions.

The current feed chain is largely downstream of the solved state,
with an optional screening combustion-efficiency coupling and an
implicit wall-temperature root. No new pump-to-nozzle feedback should
be invented merely to make every arrow two-way. Add a coupled residual
only when a physical dependency and boundary condition justify it.

\subsection{Reoptimization and fidelity reporting}
Checking an optimized TOP wall against MOC after the solve does not
produce an MOC optimum. A transition study should reoptimize using the
new accepted analysis and compare the objective, constraints and
design under flow/thermal refinement and multiple starts.

The current default dry-mass objective is explicitly partial: important
manifolds, valves, lines, gimbal/mount and other injector hardware are
not all represented. Better nozzle physics cannot complete those
component models or award hardware qualification. Keep the objective's
mass boundary, applicable constraints and model-coverage flags visible.

During migration, explicit fidelity levels can support broad searches,
candidate reanalysis and final higher-fidelity reoptimization. If a
surrogate is used, bound its discrepancy and derivative error over a
declared training/validation domain. A fallback or model switch that
changes the objective between candidates must be recorded and managed;
finite outputs alone do not make the optimization consistent.
